The common two-hour comparison contains 440 person-hours. We distinguish the full set of 24 sampled plans from the six controlled benchmarks with equal $68.2\%$ direct-service coverage. In the full set, A3 has the highest $\eta$, $37.58$ person$\cdot{}^{\circ}$C$\cdot$h/kWh. Within the equal-coverage subset, A36 is highest at $35.49$. A36 is not the efficiency maximum across all counts. These are conditional maxima over finite candidate sets, not global optima.

\begin{table}[!ht]
\centering\small
\caption{Efficiency accounting under the nominal assumptions ($\beta=0.5$, $\chi=1$). $U$ is useful relief in person$\cdot{}^{\circ}$C$\cdot$h; $E$ is kWh. Coverage describes the service weights, not a separate objective.}
\label{tab:deployment-results}
\begin{tabular}{@{}lrrrr@{}}
\toprule
Plan & Person-time coverage & $U$ & $E$ & $\eta=U/E$\\
\midrule
A3 & 5.68\% & 56.38 & 1.5 & 37.58\\
A36 & 68.18\% & 638.89 & 18.0 & 35.49\\
B24 & 68.18\% & 446.28 & 57.6 & 7.75\\
C18 & 68.18\% & 453.00 & 90.0 & 5.03\\
AB30 & 68.18\% & 581.61 & 37.8 & 15.39\\
AC27 & 68.18\% & 593.68 & 54.0 & 10.99\\
ABC24 & 68.18\% & 524.76 & 63.9 & 8.21\\
AB30-R & 76.19\% & 695.07 & 37.8 & 18.39\\
\bottomrule
\end{tabular}
\end{table}

\noindent\textbf{Scale and allocation.} Increasing the number of fans increases served person-time, while their heat also changes the background; the added useful relief need not be proportional to electricity. The resulting count sweep is evaluated through $U/E$. A3's $5.68\%$ coverage is recorded as the scope of its service, without imposing a coverage threshold in this efficiency comparison.

AB30-R redistributes the same 18 A and 12 B units from equal counts to $(3,4,0)$, $(12,3,0)$ and $(3,5,0)$. Electricity remains 37.8 kWh and area coverage $68.2\%$, while person-time coverage increases from $68.2\%$ to $76.2\%$. Useful relief rises from $581.61$ to $695.07$ person$\cdot{}^{\circ}$C$\cdot$h, increasing efficiency from $15.39$ to $18.39$. The gain results from delivering the same inventory's cooling closer to passenger demand.

\noindent\textbf{Efficiency sensitivity.} For the 15 tested combinations $\beta=0,0.25,0.5,0.75,1$ and $\chi=0,0.5,1$, A36 remains the efficiency leader among the six equal-coverage benchmarks. Changing A's circulation from $0.3$ to $0.6$ to $1.2$ kg dry air/s gives A36 efficiencies $34.06/35.49/36.22$; it retains its equal-coverage efficiency lead. The separate footprint-sensitivity figure reports efficiency changes at fixed equipment counts. These are defined parameter tests, not robustness over every possible climate or delivery pattern. Overall adequacy and tradeoffs are considered separately in Section~\ref{sec:effectiveness}.
